\chapter{Competition creates research}\label{ch:research}

\begin{chapterplan}
\textbf{The question.} When the informed investor must pay for its information, does a stronger incumbent still bring
informed trade and more entry? This chapter derives the research condition, proves Proposition~4 and reads
Corollary~3 and its figure.
\end{chapterplan}

% =====================================================================
\section{The set-up}\label{sec:research-setup}

\begin{keyidea}
Research is the presence probability chosen at a cost. More research makes a purchase better news for class \(H\), so
it reduces the camouflage of the investor's own orders. A wider spread makes research pay more, and enough research
lets a purchase recruit a high-cost challenger.
\end{keyidea}

The prior is \(1/2\). The research results take \(0<\beta<1/2\), ranked classes, a continuous incumbent law, \(\DT>0\),
\(\gamma=g_H-g_L>0\) and \(\tau>1/2\). Before it learns the class, the investor chooses \(\lambda\in[0,\lb]\), with
\(0<\lb<1\), and pays \((a/2)\lambda^2\), with \(a>0\). Research succeeds with probability \(\lambda\). After a failure the investor knows no more than the market, so it leaves
the slot to a liquidity trader. So \(b=(1-\lambda)\beta\), \(b_0=(1-\lambda)(1-2\beta)\) and \(u=\lambda/b\). Market makers and the challenger
use a conjectured intensity \(\hat\lambda\). Condition (E7), research optimality, asks that \(\hat\lambda\) be the
research choice of a best response at the fixed prices and preparation rule.

\paragraph{The posteriors.} At the prior \(1/2\) the posterior bounds of Chapter~\ref{ch:beliefs} become
\begin{equation}\label{eq:research-bounds}
m_\lambda=\frac{\beta(1-\lambda)}{2\beta+(1-2\beta)\lambda},\qquad
M_\lambda=1-m_\lambda=\frac{b+\lambda}{2b+\lambda}.
\end{equation}
The quotient rule gives \(m_\lambda'=-\beta/[2\beta+(1-2\beta)\lambda]^2<0\), so \(M_\lambda\) rises strictly, from
\(M_0=1/2\). For \(1/2<\tau<1\), the recruitment intensity \(\lambda_E(\tau)\) solves \(M_\lambda=\tau\). Solving \((1+u)/(2+u)=\tau\) for \(u\), and then for \(\lambda\),
\begin{equation}\label{eq:research-lamE}
\lambda_E(\tau)=\frac{\beta(2\tau-1)}{\beta+(1-2\beta)(1-\tau)},\qquad m_{\lambda_E(\tau)}=1-\tau .
\end{equation}
A full purchase recruits exactly when \(\lambda\ge\lambda_E(\tau)\), and \(\lambda_E\) rises with \(\tau\).

% =====================================================================
\section{The value of a success and the research condition}\label{sec:research-value}

\paragraph{Lemma A.12.} Let \(\Psi=\tfrac12\max_qU_H(q)+\tfrac12\max_qU_L(q)\) be the best trading value after a
success. A deviation does not move prices, so the best payoff of research \(\tilde\lambda\) is
\(\tilde\lambda\Psi-(a/2)\tilde\lambda^2\). This is strictly concave, so research optimality is exactly
\(\hat\lambda=\min\{\lb,\Psi/a\}\), unique, under all four tie rules.

\paragraph{The two classes.} In recruiting research the purchase recruits, so \(U_H(1)=\DT m_\lambda-k\). A sale never
recruits, so class \(L\) trades at entry \(\rho\). Its sale at rate \(s_-=\lambda\sigma_L(-1)\) earns
\(\rho\DT b/(2b+s_-)-k\), which falls strictly in \(s_-\). So (Lemma A.15) the unique sale rate is
\(s_-=\min\{\lambda,\ b(\rho\DT/k-2)^+\}\): class \(L\) abstains when \(\rho\DT\le2k\), mixes when
\(\rho\DT m_\lambda<k<\rho\DT/2\), and sells for sure when \(\rho\DT m_\lambda\ge k\). Its best payoff is
\((\rho\DT m_\lambda-k)^+\).

Average the two classes and apply Lemma A.12 below the cap. This is equation (10) of the paper:
\begin{equation}\label{eq:research-cond}
2a\lambda=\DT\,m_\lambda-k+\bigl(\rho\DT\,m_\lambda-k\bigr)^+ .
\end{equation}
The factor \(2\) appears because a success reveals class \(H\) only half the time. Let \(\mathcal R(\lambda)\) be the
right side minus the left side. The left side rises and the right side falls (less camouflage), so \(\mathcal R\) is
continuous and strictly decreasing. Since \(\mathcal R(0)\ge\tfrac12\DT-k\) and \(\mathcal R(1)=-k-2a<0\), if \(\DT>2k\) there is one root
\(\lambda_R(\DT)\in(0,1)\), and (E7) is the same as \(\hat\lambda=\min\{\lambda_R,\lb\}\). If \(\DT\le2k\), a purchase loses
money at every positive \(\lambda\).

\paragraph{The exact recruitment margin.} With \(J=\DT(1-\tau)\), the paper evaluates \(\mathcal R\) at
\(\lambda_E(\tau)\):
\begin{equation}\label{eq:research-margin}
\varphi_a^*=\mathcal R(\lambda_E(\tau))=J-k+(\rho J-k)^+-2a\lambda_E(\tau).
\end{equation}
Since \(\mathcal R\) falls and \(M_\lambda\) rises, \(M_{\min\{\lambda_R,\lb\}}\ge\tau\) holds exactly when
\(\varphi_a^*\ge0\) and \(\lambda_E(\tau)\le\lb\).

% =====================================================================
\section{Proposition 4: competition creates research}\label{sec:research-prop4}

\begin{paperresult}{Proposition 4 (Competition creates research)}
Let the classes satisfy the class ranking and the prior be \(1/2\), and let an indifferent investor mix over its best
orders and an indifferent high-cost challenger prepare. Take incumbent distributions with a positive spread,
\(g_H>g_L\) and \(\tau>1/2\), and consider the equilibrium in which class \(H\) buys after a success and a purchase
recruits a high-cost challenger. The same recruiting outcome is available when an indifferent high-cost challenger
may prepare with any probability and prepares fully at the purchase.
\begin{enumerate}[label=(\alph*),leftmargin=2em]
\item The equilibrium exists if and only if \(\DT>2k\) and \(\tau\le M_\lambda\) at \(\lambda=\min\{\lambda_R(\DT),\lb\}\),
which is then the investor's research.
\item Research rises strictly with the spread until it reaches the cap and then stays at the cap. Entry,
\(\rho+(1-\rho)[\beta+(1/2-\beta)\lambda]\), and entry in class \(H\) rise with research, and entry in class \(L\) falls.
\item When class \(H\) outbids the incumbent, high-class ownership equals half the entry in class \(H\). Along a path on
which class \(H\) outbids every incumbent, ownership does not fall with strength, and it rises strictly where research
is below the cap and the spread rises strictly.
\item Along a compact path of such distributions with continuous spread and threshold, let the spread against the
weakest incumbent lie below the value at which research reaches its cap, let \(r_{\lb}\) be the first strength at
which the spread reaches that value, and let the threshold stay below one up to \(r_{\lb}\). If the exact recruitment
margin rises strictly up to \(r_{\lb}\), from a negative value against the weakest incumbent to a positive value at
\(r_{\lb}\), then the equilibrium exists exactly on an interval of strengths. On it, research rises from the
recruitment intensity to the cap and then stays at the cap, up to the last strength at which the threshold does not
exceed the top posterior at the cap.
\end{enumerate}
\end{paperresult}

\paragraph{Proof of (a).} Positive research forces a strict full purchase, \(\mu(1)=M_{\hat\lambda}\), and entry \(\rho\)
after a sale or no order (Lemmas A.13 and A.14). So
\(\Psi=\tfrac12[\DT m_{\hat\lambda}-k+(\rho\DT m_{\hat\lambda}-k)^+]\), and (E7) gives \(\hat\lambda=\min\{\lambda_R,\lb\}\),
which needs \(\DT>2k\). The purchase recruits exactly when \(M_{\hat\lambda}\ge\tau\). Conversely, this research, full
purchases, the sale rate of Lemma A.15, entry \((\rho,\rho,1)\), competitive prices and truthful bids satisfy (E1) to
(E7).

\paragraph{Proof of (b).} \(\mathcal R\) rises strictly with \(\DT\) at each \(\lambda\), so the root \(\lambda_R(\DT)\)
rises strictly, and it is continuous. Only a purchase raises entry above \(\rho\), and a purchase has probability
\(\tfrac12(b+\hat\lambda)+\tfrac12 b=\beta+(\tfrac12-\beta)\hat\lambda\). So
\(\Ent=\rho+(1-\rho)[\beta+(\tfrac12-\beta)\hat\lambda]\). In the same way \(e_H=\rho+(1-\rho)[\beta+(1-\beta)\hat\lambda]\) and \(e_L=\rho+(1-\rho)\beta(1-\hat\lambda)\). Their
slopes in \(\hat\lambda\) are \((1-\rho)(\tfrac12-\beta)>0\), \((1-\rho)(1-\beta)>0\) and \(-(1-\rho)\beta<0\). Research adds an
informed purchase half the time and removes a liquidity purchase with probability \(\beta<1/2\).

\paragraph{Proof of (c).} Given the class, preparation is independent of the values, so \(\Own_H=\tfrac12e_HW_H(F)\),
which is \(e_H/2\) when \(W_H=1\). Along a path the spread never falls, so \(e_H\) never falls.

\paragraph{Proof of (d).} At the cap, recruitment holds exactly when \(\tau\le M_{\lb}\). The spread and the threshold
never fall, so this set is an interval that ends at \(\max\{r:\tau(r)\le M_{\lb}\}\). Below \(r_{\lb}\), recruitment is
\(\varphi_a^*\ge0\), and the rising margin has one zero \(r_R\). The paper's cheap-research condition (CR),
\[
[1+\rho\1\{\rho J>k\}]\,J'-\frac{2a\beta\,\tau'}{[\beta+(1-2\beta)(1-\tau)]^2}>0
\quad\text{almost everywhere on }[r_-,r_{\lb}],
\]
is the derivative of \eqref{eq:research-margin} along the path: a recruiting purchase's return rises faster than the
marginal cost of the research needed to reach the threshold.

% =====================================================================
\section{The three research outcomes}\label{sec:research-outcomes}

In \emph{floor-entry research}, both classes trade after a success and no order recruits. For \(\rho\DT\ge2k\), \(\lambda_Q(\DT)\) is the unique root in
\([0,1)\) of
\begin{equation}\label{eq:research-floor}
a\lambda_Q=\rho\DT\,m_{\lambda_Q}-k .
\end{equation}

Lemma A.16 of the paper shows that, under the default tie rules, three outcomes exhaust the equilibria: zero research,
exactly when \(\rho\DT\le2k\), with entry \(\rho\); recruiting research, exactly as in Proposition~4(a); and
floor-entry research, exactly when \(\rho\DT>2k\) and \(\tau>M_{\min\{\lambda_Q(\DT),\lb\}}\), with research
\(\min\{\lambda_Q,\lb\}\), full orders and entry \(\rho\). The two positive-research outcomes coexist exactly when
\(\rho\DT>2k\) and \(M_{\min\{\lambda_Q,\lb\}}<\tau\le M_{\min\{\lambda_R,\lb\}}\).

\emph{Proof sketch.} With positive research the purchase
is strict and full (Lemma A.14). If it does not recruit, both classes earn \(\rho\DT m_{\hat\lambda}-k=\Psi\), so
\(\hat\lambda=\min\{\lambda_Q,\lb\}\). At \(\lambda_Q\), \(2a\lambda_Q=2(\rho\DT m_{\lambda_Q}-k)\), so
\(\mathcal R(\lambda_Q)=(1-\rho)\DT m_{\lambda_Q}>0\) and \(\lambda_Q<\lambda_R\). Against a strong incumbent, condition
(U), \(\rho\DT m_{\lb}-k>a\lb\), with \(M_{\lb}\ne\tau\), gives a unique equilibrium under all four tie rules, with
research \(\lb\) and full orders (Lemma A.17).

% =====================================================================
\section{Corollary 3: research in the example}\label{sec:research-example}

On the example's path the spread and the threshold rise strictly. The marked strengths are \(r_R\), where
\(\varphi_a=\DT(1-\tau)-k-2a\lambda_E=0\); \(r_{\lb}\), where research reaches the cap; \(r_N\), where \(\rho\DT=2k\);
\(r_Q\), where \(\varphi_\rho=\rho\DT(1-\tau)-k-a\lambda_E=0\); and \(r_C\), where \(\tau=M_{\lb}\). Chapter~\ref{app:numbers}
lists the values.

\begin{paperresult}{Corollary 3 (Research equilibria in the example)}
In the example at the prior \(1/2\), with the paper's research cost \(a\) and cap \(\lb\), let an indifferent investor
mix over its best orders and an indifferent high-cost challenger prepare. On the path from \(r_0\) to \(r_2\):
\begin{enumerate}[label=(\roman*),leftmargin=2em]
\item zero research, with entry \(\rho\), is an equilibrium exactly on \([r_0,r_N]\);
\item recruiting research is an equilibrium exactly on \([r_R,r_C]\); research is below the cap on \([r_R,r_{\lb})\) and
at the cap on \([r_{\lb},r_C]\), and entry rises from \(\rho+(1-\rho)[\beta+(1/2-\beta)\lambda_E(r_R)]\) at \(r_R\) to
\(\rho+(1-\rho)[\beta+(1/2-\beta)\lb]\) at \(r_{\lb}\) and keeps that value up to \(r_C\);
\item floor-entry research, in which both classes trade after a success, no order recruits and entry is \(\rho\), is an
equilibrium exactly on \((r_N,r_Q)\), below the cap, and on \((r_C,r_2]\), at the cap;
\item there is no other equilibrium.
\end{enumerate}
The unique equilibrium thus has zero research at \(r_0\), recruiting research at the cap at \(r_1\) and floor-entry
research at the cap at \(r_2\); at \(r_1\) and \(r_2\) it is the equilibrium of Proposition~2 with \(\lambda=\lb\).
\end{paperresult}

\emph{Proof in outline.} Zero research holds on \([r_0,r_N]\), where \(\rho\DT\le2k\). The paper shows that
\(\varphi_a\) and \(\varphi_\rho\) rise strictly on the path, and \(\varphi_a(r_0)<0<\varphi_a(r_{\lb})\). Part (d) gives
recruiting research on \([r_R,r_C]\). Floor-entry research below the cap needs \(\varphi_\rho<0\), that is \(r<r_Q\);
above \(r_C\), Lemma A.17 applies. Lemma A.16 excludes everything else.

\begin{figure}[ht]
\centering
\includegraphics[width=0.62\linewidth]{research}
\caption{Research (panel (a)) and entry (panel (b)) along the example's path, at the prior \(1/2\), under the default
tie rules. Filled markers belong to a curve and open markers do not. Recruiting research starts at the positive level
\(\lambda_E(r_R)\). Entry first rises and later falls, as in Proposition~2.}
\label{fig:research-path}
\end{figure}

\begin{watchout}
Corollary 3 holds only in the example, with the paper's research cost and cap. Proposition~4 holds under its hypotheses.
\end{watchout}

\begin{chaptersummary}
Research is the presence probability chosen at a cost, up to a cap. In recruiting research it solves
\(2a\lambda=\DT m_\lambda-k+(\rho\DT m_\lambda-k)^+\), so it is \(\min\{\lambda_R(\DT),\lb\}\) when \(\DT>2k\), and a
purchase recruits when research is at least \(\lambda_E(\tau)\). Proposition~4 says that research, entry and class-\(H\)
entry rise with the spread, and class-\(L\) entry falls. In the example, zero research, recruiting research and
floor-entry research divide the path at \(r_R\), \(r_N\), \(r_Q\) and \(r_C\), and entry rises and later falls.
\end{chaptersummary}
